\documentclass[11pt,a4paper]{article}
\usepackage{amsmath,amssymb,amsfonts}
\usepackage{geometry}
\geometry{margin=1in}

\title{\textbf{Patent Specification: Constant-Time In-Memory Computing Architecture for Side-Channel Leakage Mitigation}}
\author{\textbf{Inventor:} K C Akshay Vishvesh}
\date{\today}

\begin{document}
\maketitle

\begin{abstract}
This invention discloses a hardware-enforced constant-time key-management and current-balancing crossbar architecture for In-Memory Computing (IMC) systems. By equalizing dynamic charge draw ($I_{total}$) across read operations regardless of stored bit representations, the system eliminates power-side-channel leakage during KV-cache token retrieval in transformer inference workloads.
\end{abstract}

\section{Background of Invention}
Prior art implementations rely on dynamic software masking or delay insertion, introducing severe performance overheads ($>30\%$). Current physical crossbars suffer from parasitic capacitive coupling, leaking secret key information during vector-matrix operations.

\section{Detailed Description \& Mathematical Model}
The proposed system enforces a constant current draw according to Equation~\ref{eq:power}:
\begin{equation}
I_{rail}(t) = I_{data}(t) + I_{dummy}(t) = K \quad (\text{Constant})
\label{eq:power}
\end{equation}
Where $I_{dummy}(t)$ is dynamically generated by a complementary crossbar differential element.

\section{Patent Claims}
\begin{enumerate}
    \item An In-Memory Computing (IMC) circuit comprising:
    a primary crossbar memory array, a dummy-charge balancing array, and a dynamic differential current regulator enforcing dynamic power invariance across distinct data read cycles.
    \item The circuit of claim 1, wherein bit-line charge equalizers operate within a clock domain boundary without software cycle delays.
\end{enumerate}

\begin{enumerate}
  \setcounter{enumi}{1}
  \item The compute-in-memory architecture of claim 1, further comprising a dynamic dummy charge-balancing circuit coupled to each bitline pair, wherein said dynamic dummy charge-balancing circuit injects a complementary differential charge equal to $Q_{dummy} = Q_{max} - Q_{actual}$ during each computational accumulation cycle to maintain constant total array current draw.

  \item The circuit of claim 2, wherein the dynamic dummy charge-balancing circuit comprises an array of programmable switch-capacitor banks controlled by an inverted weight-activation decode matrix operating in parallel with wordline activation.

  \item The compute-in-memory architecture of claim 1, wherein the hardware-enforced constant-time execution pipeline guarantees uniform memory-access latency and Hamming weight power dissipation across all dynamic weight loading operations within the Key-Value cache.

  \item The circuit of claim 2, further comprising a localized sense-amplifier power-rail regulator configured to damp differential voltage spikes induced by dummy charge injection within less than 150 picoseconds.
\end{enumerate}

\end{document}
